\chapter{Discusión}

% --------------------------------------------------------
\section{¿El ML supera al ARIMA?}
% --------------------------------------------------------
% TODO: ~400 palabras
% - Sí, sustancialmente: 49-64 pb vs 107.48 pb (40-54% reducción)
% - El salto validación→test del ARIMA revela que no captura shocks exógenos
% - El XGBoost captura la persistencia AR + shocks contemporáneos del mercado
% - Relevancia económica: 1 pb EMBI ≈ X millones USD en costo de deuda

% --------------------------------------------------------
\section{¿El NLP agrega valor predictivo?}
% --------------------------------------------------------
% TODO: ~400 palabras
% - Sí: 2.30-2.59 pb de mejora incremental directa
% - El NLP que sobrevive feature selection es volumen (event_count, total_articles)
%   no sentimiento (avg_tone, goldstein eliminados)
% - Interpretación: el mercado reacciona a que "se hable de Ecuador" independientemente
%   del tono — incertidumbre mediática como proxy de riesgo
% - Granger negativo: la cobertura no anticipa, es contemporánea
% - Implicación: la relación NLP-EMBI no genera arbitraje temporal

% --------------------------------------------------------
\section{¿Es el EMBI ecuatoriano exógeno?}
% --------------------------------------------------------
% TODO: ~500 palabras — la sección más rica de la discusión
% - En períodos normales: sí (Pre-COVID: top features son EMBI lag + WTI + Treasury + ETF)
% - En crisis (COVID 2020): el desempleo local irrumpe como segundo factor (SHAP=125.98)
%   → el mercado evaluaba capacidad de pago local ante el default inminente
% - Post-COVID (2022-2025): Fed + DXY dominan absolutamente
%   → coherente con la dolarización: Ecuador no tiene política monetaria propia
% - Conectar con literatura: hipótesis de "sudden stops", Calvo et al.

% --------------------------------------------------------
\section{Contexto histórico: el default Correa 2008}
% --------------------------------------------------------
% TODO: ~300 palabras
% - En 2008-2009 Ecuador repudió ~$3.2B en deuda bajo Correa (bonos 2012 y 2030)
% - Decisión política, no fiscal — primer default voluntario moderno en AL
% - El EMBI llegó a 4.000+ pb durante el evento
% - Por qué no está modelado: GDELT v2 empieza en 2013; incluir GDELT v1 requeriría
%   reconciliar dos versiones con dimensiones distintas → trabajo futuro
% - Relevancia narrativa: confirma que Ecuador tiene régimen de riesgo propio
%   diferente al promedio EM

% --------------------------------------------------------
\section{Sobre la representación NLP: ¿lineal o no lineal?}
% --------------------------------------------------------
% TODO: ~300 palabras
% - Autoencoder no mejora sobre PCA (3.3 pb diferencia)
% - PCA captura 77.7% de varianza NLP en 3 componentes → estructura esencialmente lineal
% - Las 4 variables originales (rolling 30d) ya son una representación casi óptima
% - Implicación: para mejorar la señal NLP se necesitan datos de texto crudo
%   (embeddings BERT/FinBERT sobre noticias GDELT), no compresión de métricas agregadas

% --------------------------------------------------------
\section{Limitaciones}
% --------------------------------------------------------
% TODO: ~400 palabras
% 1. Dato diario vs frecuencia de publicación de variables macro (mensual/trimestral)
%    → solución: lags y rolling means, pero se pierde precisión intra-mes
% 2. GDELT como proxy mediático: no captura texto completo, solo eventos codificados
% 3. El período post-COVID (2024-2025) es atípico: reducción de EMBI 1.800→500 pb
%    bajo Noboa — posible "nuevo régimen" que el modelo histórico no captura bien
% 4. Ausencia de datos pre-2013 (GDELT v2) excluye el default 2008
% 5. El test de Wilcoxon con 5 folds tiene bajo poder estadístico
